\section*{Prefacio}
\addcontentsline{toc}{section}{Prefacio}

El presente documento es el Plan de Gestión del Proyecto (SPMP) del Kit Pagos Colombia, un framework de software construido como trabajo de grado en Ingeniería de Sistemas de la Pontificia Universidad Javeriana. El proyecto es desarrollado por cuatro estudiantes bajo la dirección del Ing. Jaime Andrés Pavlich Mariscal, con el propósito de producir un artefacto que simplifique la integración de las principales pasarelas de pago colombianas en aplicaciones de software.

El framework tiene cuatro componentes. El SDK de servidor ofrece una interfaz unificada para Wompi, Rapyd (que adquirió PayU GPO), Mercado Pago y Kushki, y el módulo de navegador tokeniza la tarjeta en la página del comprador. La API de Simulación, desplegada en la nube, permite probar sin credenciales reales y expone también las operaciones del kit por REST. El cuarto componente es una documentación centralizada de datos de prueba adaptados al contexto colombiano. El impacto esperado es reducir la complejidad y el tiempo que los equipos de desarrollo invierten al integrar múltiples proveedores de pago, concentrando en un único punto de entrada la lógica que hoy está distribuida entre cuatro fuentes distintas.

Este documento describe cómo se ejecuta el proyecto: las decisiones metodológicas que lo estructuran, la organización del equipo, el plan de trabajo, los mecanismos de seguimiento y los procesos de calidad de cada fase. No describe el producto en detalle técnico; para eso están el SRS y el SAD.

La audiencia principal de este SPMP es el equipo de desarrollo, que lo utiliza como referencia operativa. También es relevante para el director del trabajo de grado en su rol de supervisión y aprobación de entregables, y para los evaluadores académicos que revisen el proceso de ejecución. Puede ser útil, adicionalmente, para otros estudiantes que desarrollen proyectos similares bajo Design Science Research.
